\begin{textsample}{POS Dim 3 – sp – Score 10.00 – sp/sp\_inf\_papel\_do\_professor.txt}  \label{ex:f3_pos_198}
Papel dos profissionais da instituição de Educação \textbf{Infantil} A instituição de Educação \textbf{Infantil} está centrada no atendimento aos \textbf{bebês} e às \textbf{crianças}, que estão sob a responsabilidade dos \textbf{adultos} com os quais estabelece vínculos estáveis e seguros, como os professores e cuidadores, bem como daqueles com os quais interagem ao longo da \textbf{rotina}, como os responsáveis pela \textbf{limpeza}, pela alimentação, pela segurança, pela secretaria, pela gestão da instituição, dentre outros.
Nesse sentido, é essencial que todos os profissionais conheçam as especificidades da faixa etária atendida, a fim de compreender a importância de suas ações em favor da \textbf{criança}, de modo a zelar e contribuir efetivamente com a qualidade do atendimento prestado.
Assim, também, é relevante cuidar das narrativas por meio dos quais nos dirigimos às \textbf{crianças}, nas diferentes situações do cotidiano, compreendendo esses momentos como referências de práticas sociais, que precisam ser apresentados de modo ético e empático, cientes de que as \textbf{crianças} aprendem não apenas pelo que lhe falamos, mas que observam, replicam e reinventam o que fazemos.
Por fim, é importante ressaltar que todos os profissionais que atuam direta ou indiretamente na Educação \textbf{Infantil}, assim como nas demais etapas da Educação Básica, que de algum modo participam do processo aprendizagem e desenvolvimento da \textbf{criança}, ou que deem suporte pedagógico, tornam -se corresponsáveis pela formação integral da \textbf{criança}, sendo assim considerados educadores.
Papel do professor de Educação \textbf{Infantil} Os professores da Educação \textbf{Infantil} devem priorizar o protagonismo da \textbf{criança}.
Para tanto, precisam praticar a \textbf{escuta} ativa e a mediação do processo de aprendizagem e desenvolvimento, fazendo com que as ações do cotidiano e do imaginário ( faz de conta ) se abram, intencionalmente, como um mapa de possibilidades educacionais, criando oportunidades, situações, propondo experiências que ampliem os horizontes culturais, artísticos, científicos e tecnológicos das \textbf{crianças}.
Dessa forma, é preciso compreender seu papel fundamental no desenvolvimento das \textbf{crianças} : sua intencionalidade educativa se expressa nas atividades propostas e na gestão de ambientes que promovam as interações e a \textbf{brincadeira}.
Para realizar plenamente o trabalho como professor de Educação \textbf{Infantil}, é imprescindível aprender a interpretar os processos contínuos e compreender as percepções, as ideias e os pensamentos das \textbf{crianças} sobre as ações dos \textbf{adultos} e de seus pares.
Assim, os professores devem estar atentos e conscientes sobre os interesses que surgem no decorrer do dia, durante as \textbf{brincadeiras}, e saber correlacioná -los aos objetivos de aprendizagem, conferindo sentido pedagógico às suas próprias intervenções.
Os professores devem também conhecer as bases científicas do desenvolvimento da \textbf{criança} nas diferentes fases, de \textbf{bebês} a \textbf{crianças} \textbf{pequenas}, compreendendo que as ações de educar e cuidar são práticas indissociáveis.
Também é importante garantir aos professores que continuem seu processo de aperfeiçoamento, de forma a ir além da formação inicial, assegurando formação continuada em seus espaços de trabalho, a fim de potencializar a reflexão sobre as práticas pedagógicas e construir um olhar \textbf{criterioso} sobre a aprendizagem e o desenvolvimento das \textbf{crianças}.
Para tanto, os professores precisam ser pesquisadores das práticas pedagógicas, compreendendo a necessidade de planejar com base no conhecimento específico sobre cada faixa etária, garantindo os objetivos de aprendizagem e desenvolvimento e organizando os tempos, espaços e materiais adequados à cada proposta.
Para que os objetivos sejam atingidos, os professores necessitam ser exímios observadores e fazer diferentes \textbf{registros} sobre o que observam.
É o que pode dar sustentação às avaliações, à reflexão sobre a aprendizagem e, então, às propostas para (re)encaminhamentos que garantam aprofundamento no domínio das competências e habilidades previstas para a fase.
Por fim, é importante compreender como se dá essa relação do cuidar e educar, considerada imprescindível para a construção dos saberes, a constituição do sujeito, a aprendizagem e o desenvolvimento, cientes de que o espaço e o tempo vividos pela \textbf{criança} demandam intervenções responsivas dos professores, que devem planejar vivências e ampliar as experiências a partir dos interesses e das necessidades das \textbf{crianças}.

% matched lemmas: adulto, bebê, brincadeira, criança, criterioso, escuta, infantil, limpeza, pequeno, registro, rotina
\end{textsample}
